% ===========================================================================
%  sviluppatore/06-dashboard/06-stream-vuoti.tex
% ===========================================================================
\section{Gli stream vuoti}\label{sez:stream-vuoti}

Uno stream nasce dai risultati di un backtest (\vediSez{data-layer}), un
progetto si compone di stream, un runner del Live si costruisce su un
progetto. La catena ha una conseguenza: senza backtest una strategia non entra
in un progetto, quindi non arriva al Live se non da sola. Per una prova dal
vivo --- o per vedere girare una strategia prima di averla misurata --- la
scheda di \voce{Run backtest} crea uno stream \emph{vuoto}: c'è, si mette nei
progetti, anche in progetti fatti solo di stream vuoti, e non ha metriche.

\begin{perche}{Perché il Live non ha bisogno dello stream}
Il runner fa girare il \emph{codice} delle strategie del progetto. Dello
stream usa due cose, entrambe facoltative: la vol di backtest degli approcci
per il warmup della normalizzazione (\simbolo{\_vol\_di\_backtest}, che salta
chi non ha dati e lascia lo scale a 1) e niente altro: i pesi dal vivo li
stima \file{live/ensemble\_live.py}, che parte con pesi uguali finché non ha
abbastanza osservazioni sue. Un progetto di soli stream vuoti, su TEST, parte
e cicla come un altro: lo verifica la prova nel browser.
\end{perche}

\subsection{Il segnaposto}

\file{dashboard/stream\_vuoti.py} scrive
\file{output/strategies/<nome>\_empty\_stream.json}: tipo, nome, quando. Il
posto non è casuale: sta fra i risultati perché ne segue la vita senza codice
in più.

\begin{itemize}
  \item \file{stream\_cleanup} conosce il suffisso \cod{\_empty\_stream}:
        eliminando la strategia (\voce{Delete} della scheda) il segnaposto se
        ne va con i risultati, e uno rimasto senza strategia è un residuo che
        la Home propone di togliere.
  \item Nascondere la strategia rinomina la sua cartella e lascia il
        segnaposto dov'è: lo stream vuoto sparisce e torna con lei.
  \item Il backtest lo rende inutile: \simbolo{DataLayer.startup}, dopo aver
        caricato gli stream, toglie i segnaposto di quelli che ora hanno dati
        (\simbolo{pulisci\_superati}). Da lì lo stream è uno stream e basta.
\end{itemize}

\subsection{Chi è vuoto: calcolato, non tenuto}

\simbolo{DataLayer.get\_empty\_streams} dice quali stream sono vuoti a ogni
richiesta: segnaposto presente, strategia o approccio \emph{in vista},
stream registrato (la strategia dichiara \cod{output\_path\_pattern}: senza,
non può avere uno stream, nemmeno vuoto), nessun dato.
\simbolo{stream\_state(nome)} riassume in tre parole: \cod{data},
\cod{empty}, \cod{missing}.

\begin{attenzione}{Perché non è un attributo del DataLayer}
Le ricariche dei dati costruiscono un DataLayer nuovo e ne \emph{copiano} gli
attributi in quello che le rotte hanno catturato --- a mano, in più punti di
\file{app.py}. Un attributo in più sarebbe quello che una delle copie
dimentica, e la pagina mostrerebbe gli stream vuoti di prima. Calcolarlo dal
registro e dal disco evita il problema; una piccola cache lo rilegge solo se
cambiano il registro o le cartelle \file{strategies/}, \file{approaches/},
\file{output/strategies/}. Per chi ricarica, c'è anche
\simbolo{DataLayer.adotta(nuovo)}, che fa la copia in un posto solo.
\end{attenzione}

\subsection{Metriche: solo di chi ha dati}

\simbolo{compute\_custom\_ensemble} divide i nomi in tre gruppi:

\begin{center}
\begin{tabular}{@{}L{34mm}L{100mm}@{}}
\toprule
\textbf{Gli stream} & \textbf{La risposta} \\
\midrule
\rigaTenue
nessuno con dati & \cod{no\_metrics: true}, con \cod{empty\_streams} e
                   \cod{missing\_streams}: la pagina mostra la spiegazione e
                   il progetto si salva lo stesso, con metriche vuote \\
uno con dati     & le sue metriche, con \cod{equal}: con una serie ogni modo
                   di pesare dà il \num{100}\%, e \cod{equal} non può fallire
                   (HRP e la minima varianza vogliono due serie). La risposta
                   lo dice con \cod{weighting\_mode\_requested} \\
\rigaTenue
due o più        & l'ensemble di quelli, come sempre, con l'elenco di chi
                   resta fuori \\
\bottomrule
\end{tabular}
\end{center}

\percorsoHTTP{/api/ensemble\_custom} rifiuta solo la richiesta in cui
\emph{nessun} nome è uno stream noto. Prima rifiutava anche un progetto
salvato con uno stream che non c'era più (risultati eliminati, strategia
nascosta): ora si apre, calcolato su chi c'è, con la nota.

\subsection{Il marketplace}

\simbolo{rifiuto\_pubblicazione(kind, nome, stato\_di, progetto)} è l'unico
punto della regola: uno stream vuoto non si pubblica (\cod{STRATEGY},
\cod{APPROACH}), un progetto che contiene uno stream vuoto --- o uno stream
senza dati --- nemmeno. La chiamano la rotta di pubblicazione, subito dopo aver
trovato l'artefatto e prima di ogni altro controllo (\cod{409} con il motivo
in \cod{error}), e la porta della pubblicazione automatica
(\simbolo{\_pubblica\_gratis}); \file{autopubblica.py} registra
\cod{empty\_stream} ed \cod{empty\_streams} come conclusioni \emph{non
definitive}, che si riesaminano quando arrivano i backtest. Il messaggio nomina
gli stream:

\begin{quote}\small\sffamily
A project cannot be made of empty streams: in «Live test», «live\_test\_rules»
and «live\_test\_model» are empty streams, with no backtest results. Run their
backtests, then publish the project.
\end{quote}

\subsection{Rotte e pagine}

Le rotte sono in \file{dashboard/ui\_ext\_stream\_vuoti.py}:
\rotta{GET}{/api/streams/empty},
\rotta{GET, POST}{/api/strategies/<tipo>/<nome>/empty-stream} (lo
stato, con i progetti che lo usano, oppure la creazione) e
\rotta{POST}{.../empty-stream/remove}. Gli stream vuoti arrivano anche,
\emph{a parte}, in \percorsoHTTP{/api/streams} e \percorsoHTTP{/api/a/overview}
(chiave \cod{empty\_streams}: chi legge le metriche degli stream non se li
trova mescolati), nello stato degli stream di ogni progetto in
\percorsoHTTP{GET /api/projects} (calcolato, mai salvato), in
\percorsoHTTP{/api/strategies} (\cod{stream}, \cod{stream\_possible} di ogni
scheda) e in \percorsoHTTP{/api/stream/<nome>}. La pagina di uno stream vuoto è
un ramo di \file{stream.html}: niente pannelli di metriche, una spiegazione.

\begin{pratica}{Provarli}
\file{tests/prova\_stream\_vuoti.py} li prova su una copia del progetto, con
strategie scritte dalla prova e risultati veri in \file{output/}: le schede, i
rifiuti, gli elenchi, l'ensemble nei tre casi, i progetti, il marketplace a
mano e automatico, il Live, gli strumenti dell'agente, il backtest che
sostituisce il segnaposto, la vita del segnaposto. \num{75} controlli.
\end{pratica}
